\documentclass{ctexart}

% 支持从项目根目录或当前目录编译。
\makeatletter
\def\input@path{{../}}
\makeatother
\usepackage{researchnotes}

\title{极限定义习题}
\author{}
\date{}

\begin{document}
\maketitle

按极限定义（$\varepsilon$--$\delta$ 法）证明：
\[
  \lim_{x\to 1}\sqrt{\frac{7}{16x^2-9}}=1.
\]

\section{分析：怎样找到合适的邻域}

\subsection{明确要控制的误差}

按照极限的 $\varepsilon$--$\delta$ 定义，必须证明：对任意给定的
$\varepsilon>0$，都存在 $\delta>0$，使得定义域内的 $x$ 只要满足
$0<|x-1|<\delta$，就有
\[
  \left|\sqrt{\frac{7}{16x^2-9}}-1\right|<\varepsilon.
\]
这里，$\varepsilon$ 表示对函数值误差的要求，它是任意给定的，不能只选取
某个方便的正数；$\delta$ 表示为满足这一要求而对 $x$ 与 $1$ 的距离施加的限制，
需要根据 $\varepsilon$ 来确定，而且不能依赖具体的 $x$。

因此，我们先不直接写出 $\delta$，而是从上述误差不等式出发，寻找能够保证它成立的条件。
下面的分析用于寻找 $\delta$，随后再按定义的顺序给出正式证明。

\subsection{有理化：把误差与 \texorpdfstring{$|x-1|$}{|x-1|} 联系起来}

已知条件控制的是 $|x-1|$，而目标中出现的是根式与 $1$ 的差。
为了把两者联系起来，利用平方差公式进行有理化。只要 $16x^2-9>0$，便有
\begin{align*}
  \left|\sqrt{\frac{7}{16x^2-9}}-1\right|
  &=\frac{\left|\frac{7}{16x^2-9}-1\right|}
          {\sqrt{\frac{7}{16x^2-9}}+1}\\
  &=\frac{|7-(16x^2-9)|}
          {(16x^2-9)\left(\sqrt{\frac{7}{16x^2-9}}+1\right)}\\
  &=\frac{16|1-x^2|}
          {(16x^2-9)\left(\sqrt{\frac{7}{16x^2-9}}+1\right)}\\
  &=\frac{16|x+1|}
          {(16x^2-9)\left(\sqrt{\frac{7}{16x^2-9}}+1\right)}\,|x-1|.
\end{align*}
最后一步使用了 $1-x^2=(1-x)(1+x)$。现在误差中已经出现了 $|x-1|$，
但它前面的系数仍随 $x$ 变化。接下来的目标是在 $1$ 的某个固定邻域内，
用一个与 $x$ 无关的正常数控制这个系数。

\subsection{先选固定邻域，控制其余因子}

根式的定义域由 $16x^2-9>0$ 决定，即
\[
  x<-\frac34\quad\text{或}\quad x>\frac34.
\]
在 $1$ 附近，首先需要避免 $x$ 接近 $\frac34$，否则分母可能过小。
$1$ 到 $\frac34$ 的距离是 $\frac14$，我们取这个距离的一半作为初步的邻域半径，
要求
\[
  |x-1|<\frac18.
\]
这就给出 $\frac78<x<\frac98$，从而
\[
  16x^2-9>16\left(\frac78\right)^2-9=\frac{13}{4}>3,
  \qquad
  |x+1|=x+1<\frac{17}{8}<3.
\]
此外，根式为正，所以
\[
  \sqrt{\frac{7}{16x^2-9}}+1>1.
\]
对分子使用上界、对分母使用正的下界，便得到
\[
  \frac{16|x+1|}
       {(16x^2-9)\left(\sqrt{\frac{7}{16x^2-9}}+1\right)}
  <\frac{16\cdot3}{3\cdot1}=16.
\]
因此，当 $0<|x-1|<\frac18$ 时，
\begin{equation}\label{eq:limit-error-bound}
  \left|\sqrt{\frac{7}{16x^2-9}}-1\right|<16|x-1|.
\end{equation}
这里使用 $3$ 和 $1$ 是为了简化估计，不必求出最精确的界。
同样，$\frac18$ 也不是唯一的选择；它的作用是提供一个分母远离零的固定邻域。

\subsection{根据误差要求反推 \texorpdfstring{$\delta$}{delta}}

由式~\eqref{eq:limit-error-bound}，要使函数值误差小于 $\varepsilon$，
只需进一步要求
\[
  16|x-1|<\varepsilon,
  \qquad\text{即}\qquad
  |x-1|<\frac{\varepsilon}{16}.
\]
这是一个充分条件，不要求它与原来的误差不等式等价。
不过，前面的估计是在 $|x-1|<\frac18$ 的条件下得到的，因此必须同时保证
\[
  |x-1|<\frac18
  \quad\text{和}\quad
  |x-1|<\frac{\varepsilon}{16}.
\]
于是取这两个上界中较小的一个：
\[
  \delta=\min\left\{\frac18,\frac{\varepsilon}{16}\right\}.
\]
这样，$|x-1|<\delta$ 就同时满足两个要求。由于 $\varepsilon>0$，这个 $\delta$
一定为正，且只依赖 $\varepsilon$。不能在未经说明的情况下只取
$\delta=\frac{\varepsilon}{16}$：当 $\varepsilon$ 较大时，这个半径可能超过
$\frac18$，便不能直接使用已经建立的估计。

\section{按极限定义完成证明}

\begin{proof}

任给 $\varepsilon>0$，取
\[
  \delta=\min\left\{\frac18,\frac{\varepsilon}{16}\right\}.
\]
当 $0<|x-1|<\delta$ 时，由 $|x-1|<\frac18$ 得
\[
  \frac78<x<\frac98,
  \qquad
  16x^2-9>16\left(\frac78\right)^2-9=\frac{13}{4}>3.
\]
因此根式有定义，且 $|x+1|<3$。有理化后可得
\begin{align*}
  \left|\sqrt{\frac{7}{16x^2-9}}-1\right|
  &=\frac{\left|\frac{7}{16x^2-9}-1\right|}
          {\sqrt{\frac{7}{16x^2-9}}+1}\\
  &=\frac{16|x-1|\,|x+1|}
          {(16x^2-9)\left(\sqrt{\frac{7}{16x^2-9}}+1\right)}\\
  &<\frac{16|x-1|\cdot3}{3\cdot1}
   =16|x-1|
   <16\delta
   \leq\varepsilon.
\end{align*}
这说明对任意 $\varepsilon>0$，上述 $\delta>0$ 都满足极限定义的要求，故
\[
  \lim_{x\to1}\sqrt{\frac{7}{16x^2-9}}=1.
\]
\end{proof}

\end{document}
